% Part: many-valued-logic
% Chapter: sequent-calculus
% Section: introduction

\documentclass[../../../include/open-logic-section]{subfiles}

\begin{document}

\olfileid{mvl}{seq}{int}

\olsection{પરિચય}

શાસ્ત્રીય તર્કશાસ્ત્રનું સિક્વન્ટ કલન કાર્યક્ષમ અને સરળ
!!{derivation} તંત્ર છે. જો કોઈ બહુમૂલ્ય તર્કશાસ્ત્ર સાન્ત
સત્યમૂલ્યો ધરાવતી શ્રેણિક વડે વ્યાખ્યાયિત થાય, એટલે કે $V$
સાન્ત હોય, તો તેના માટે સિક્વન્ટ કલન આપી શકાય છે.
તે કેવી રીતે કરવું તેનો વિચાર શાસ્ત્રીય કિસ્સામાં સિક્વન્ટોના
અર્થ અને અનુમાનના નિયમોના સ્વરૂપ પરથી આવે છે.

હવે યાદ કરો કે સિક્વન્ટ
\begin{align*}
    !A_1, \dots, !A_m & \Sequent !B_1, \dots, !B_n
\intertext{નું અર્થઘટન નીચેના !!{formula} તરીકે થઈ શકે છે:}
    (!A_1 \land \cdots \land !A_m) & \lif (!B_1 \lor \cdots \lor
!B_n)
\end{align*}
\sourcecorrection{OLMV-017}{સ્થિર અંગ્રેજી મૂળ
સિક્વન્ટની ડાબી બાજુનો છેલ્લો ક્રમસૂચક $n$ લખે છે,
પણ તરત નીચેના અર્થઘટનમાં ડાબી બાજુ માટે $m$ વાપરે છે.
બંને બાજુની લંબાઈ સ્વતંત્ર રહે તે માટે અહીં ડાબે
$!A_1,\dots,!A_m$ અને જમણે $!B_1,\dots,!B_n$ રાખ્યા છે.}
બીજા શબ્દોમાં, !!^a{valuation}~$\pAssign{v}$ સિક્વન્ટ
$\Gamma \Sequent \Delta$ને \emph{સંતોષે છે} જો અને માત્ર જો
કોઈ $!A \in \Gamma$ માટે $\pValue{v}(!A) = \False$ હોય
અથવા કોઈ $!A \in \Delta$ માટે $\pValue{v}(!A) = \True$ હોય.
આ અર્થઘટન મુજબ આરંભિક સિક્વન્ટો $!A \Sequent !A$
હંમેશા સંતોષાય છે, કારણ કે $\pValue v(!A) = \True$ અથવા
$\pValue{v}(!A) = \False$ હોય છે.
\sourcecorrection{OLMV-018}{સ્થિર અંગ્રેજી મૂળ
આરંભિક સિક્વન્ટ અંગેના અસત્ય-કિસ્સામાં
\texttt{\textbackslash pValue}નું મૂલ્યાંકન-ચલ $v$
છોડી દે છે. બંને વિકલ્પોમાં એ જ $v$ રાખ્યું છે.}

અહીં બાજુનાં સૂત્રો $\Gamma$ અને $\Delta$ છોડી દઈને
$\Log{LK}$માં પ્રેરણ સંયોજકના અનુમાન-નિયમો આપ્યાં છે:

\begin{defish}
    \Axiom$ \fCenter !A$
    \Axiom$ !B \fCenter $
    \RightLabel{\LeftR{\lif}}
    \BinaryInf$ !A \lif !B \fCenter $
    \DisplayProof
    \hfill
    \Axiom$ !A \fCenter !B$
    \RightLabel{\RightR{\lif}}
    \UnaryInf$ \fCenter !A \lif !B $
    \DisplayProof
\end{defish}

ઉપર આપેલા સિક્વન્ટના અર્થઘટનને લાગુ કરીએ તો
$\LeftR{\lif}$ નિયમ કહે છે કે જો $\pValue v(!A) = \True$
અને $\pValue v(!B) = \False$ હોય, તો
$\pValue v(!A \lif !B) = \False$ હોય છે. તેવી જ રીતે,
$\RightR{\lif}$ નિયમ કહે છે કે જો $\pValue v(!A) = \False$
અથવા $\pValue v(!B) = \True$ હોય, તો
$\pValue v(!A \lif !B) = \True$ હોય છે. હકીકતમાં,
આ બંને શરતી વિધાનો દ્વિશરતી છે. $\LeftR{\land}$ અને
$\RightR{\lor}$ નિયમોના પ્રમાણભૂત સ્વરૂપમાં તેમના
અનુરૂપ શરતી વિધાનો દ્વિશરતી નહીં હોય. પરંતુ આ
નિયમોનાં વૈકલ્પિક સ્વરૂપો છે જેમાં તે દ્વિશરતી બને છે:

\begin{defish}
    \Axiom$!A, !B, \Gamma \fCenter \Delta$
    \RightLabel{\LeftR{\land}}
    \UnaryInf$!A \land !B, \Gamma \fCenter \Delta$
    \DisplayProof
    \hfill
    \Axiom$ \Gamma \fCenter \Delta, !A, !B$
    \RightLabel{\RightR{\lor}}
    \UnaryInf$ \Gamma \fCenter \Delta, !A \lor !B$
    \DisplayProof
\end{defish}

આ મૂળભૂત વિચાર $n$-મૂલ્ય તર્કશાસ્ત્ર પર લાગુ કરતાં
બેને બદલે $n$ સ્થાનોવાળું સિક્વન્ટ કલન મળે છે, જેમાં
દરેક સત્યમૂલ્ય માટે એક સ્થાન હોય છે. $V = \{\False, \Undef, \True\}$
ધરાવતા ત્રિમૂલ્ય તર્કશાસ્ત્ર માટે સિક્વન્ટ
$\Gamma \mid \Pi \mid \Delta$ સ્વરૂપનું પદ છે. તે
!!a{valuation}~$\pAssign v$માં સંતોષાય છે જો અને માત્ર જો
કોઈ $!A \in \Gamma$ માટે $\pValue{v}(!A) = \False$
અથવા કોઈ $!A \in \Delta$ માટે $\pValue{v}(!A) = \True$
અથવા કોઈ $!A \in \Pi$ માટે $\pValue{v}(!A) = \Undef$
હોય. તેથી આરંભિક સિક્વન્ટો $!A \mid !A \mid !A$
હંમેશા સંતોષાય છે.

\end{document}
